\chapter{Planificación del proyecto}
\label{chap:planificacion}

La planificación del proyecto toma como referencia las prácticas de
dirección de proyectos del Project Management Institute (PMI). Se han
adaptado al alcance de este TFG instrumentos como el acta de constitución,
la estructura de descomposición del trabajo (EDT) y el seguimiento de
plazos, riesgos y cambios \cite{pmi2021pmbok}.

\section{Acta de constitución del proyecto}
\label{sec:acta_constitucion}

La Tabla~\ref{tab:acta_constitucion} recoge el propósito, el alcance,
las restricciones y el criterio de cierre definidos para el proyecto.

\begin{longtable}{@{}p{4cm}p{\dimexpr\textwidth-4cm-2\tabcolsep\relax}@{}}
\toprule
\textbf{Elemento} & \textbf{Descripción} \\
\midrule
\endfirsthead

\multicolumn{2}{@{}l}{\textit{Acta de constitución del proyecto (continuación)}}\\
\toprule
\textbf{Elemento} & \textbf{Descripción} \\
\midrule
\endhead

Propósito &
Desarrollar una herramienta que permita generar automáticamente modelos de características en formato UVL para facilitar la creación de modelos de variabilidad software, proporcionando a investigadores y desarrolladores un mecanismo configurable para definir las características, restricciones y atributos de los modelos generados. \\

\midrule

Objetivos &
Los objetivos del proyecto son los definidos en la Sección~\ref{sec:objetivos}. De forma general, el proyecto busca diseñar e implementar un generador automático de modelos de características, integrarlo en UVLHub y proporcionar soporte para diferentes niveles de expresividad del lenguaje UVL. \\

\midrule

Alcance &
El proyecto comprende el diseño e implementación del generador automático de modelos de características, su integración dentro de la plataforma UVLHub y la incorporación progresiva de soporte para modelos booleanos, aritméticos y con tipos definidos por UVL. Quedan fuera del alcance la modificación del funcionamiento interno de Flamapy y el desarrollo de nuevos algoritmos de análisis o razonamiento. \\

\midrule

Supuestos &
Se asume la disponibilidad del ecosistema software necesario para el desarrollo, incluyendo acceso al código fuente de UVLHub y Flamapy, así como las herramientas necesarias para la generación y análisis de modelos de variabilidad. También se considera disponible el entorno de desarrollo necesario para implementar, probar y validar la solución. \\

\midrule

Restricciones &
El proyecto está condicionado por las tecnologías y arquitectura existentes en UVLHub y Flamapy, especialmente el lenguaje Python (su \textit{framework} Flask) y el estándar UVL. Además, debe desarrollarse antes del 7 de octubre, la fecha límite para la entrega de la memoria. \\

\midrule

Entregables principales &
Los principales entregables del proyecto son el código fuente del generador automático integrado en UVLHub y la memoria del TFG. \\

\midrule

Criterio de cierre &
El proyecto se considerará finalizado cuando el generador produzca
modelos válidos para los niveles de expresividad incluidos en el alcance,
se encuentre integrado en UVLHub y se hayan ejecutado pruebas de la
generación y de su integración con el asistente existente. \\

\bottomrule
\caption{Acta de constitución del proyecto.}
\label{tab:acta_constitucion}
\end{longtable}

\subsection{Interesados del proyecto}
\label{sec:interesados}

La Tabla~\ref{tab:interesados} identifica a los participantes en el
proyecto y resume su papel en el desarrollo o la evaluación del TFG.

\begin{longtable}{@{}P{3cm} P{3.2cm} P{5.5cm} P{2.5cm}@{}}
  \toprule
  \textbf{Interesado} & \textbf{Rol} & \textbf{Interés / expectativa} & \textbf{Implicación} \\
  \midrule
  \endhead

  Autor del proyecto &
  Responsable del desarrollo &
  Completar el Trabajo Fin de Grado mediante el desarrollo de una solución funcional, documentada y alineada con los objetivos establecidos. &
  Diaria \\[0.3em]

  Tutores del TFG &
  Supervisión y validación académica &
  Garantizar el correcto desarrollo del proyecto, proporcionando orientación técnica y revisando los avances y entregables realizados. &
  Periódica \\[0.3em]

  Grupo de investigación / ecosistema UVLHub-Flamapy &
  Entorno receptor de la solución &
  Disponer de una herramienta que permita generar modelos de características de forma automática y ampliar las capacidades existentes del ecosistema de variabilidad software. &
  Puntual \\[0.3em]

  Tribunal del TFG &
  Evaluación final del proyecto &
  Valorar la calidad de la memoria, la defensa realizada y los resultados obtenidos durante el desarrollo del trabajo. &
  Al cierre \\

  \bottomrule
  \caption{Registro de interesados del proyecto.}
  \label{tab:interesados}
\end{longtable}

\section{Metodología}
\label{sec:metodologia}

\subsection{Metodología de desarrollo}
\label{sec:metodologia_desarrollo}


Para el desarrollo del proyecto se ha seguido una metodología incremental e iterativa, adaptada a la naturaleza evolutiva del generador automático de modelos de características. Debido al carácter investigador y evolutivo del proyecto, la especificación inicial se refinó progresivamente durante el desarrollo conforme se analizaban las necesidades del generador y su integración dentro del ecosistema existente.

En una primera etapa se desarrolló una versión inicial del generador orientada a la creación de modelos de características booleanos, junto con su integración dentro de UVLHub mediante la interfaz gráfica proporcionada para configurar el proceso de generación. Una vez obtenida esta primera versión funcional, se realizaron ciclos sucesivos de implementación, validación y corrección que permitieron mejorar la estabilidad del sistema e incorporar nuevas funcionalidades pertinentes a los siguientes niveles de expresividad que se pretendió conseguir. Esta forma de hacerlo permitió validar cada nueva incorporación antes de continuar con las siguientes etapas y reducir el riesgo de introducir errores en funcionalidades previamente implementadas que se pudieran arrastrar a las siguientes.


\subsubsection*{Alternativas consideradas}

No se optó por metodologías secuenciales como el modelo en cascada, ya que la naturaleza del proyecto requería realizar ajustes continuos conforme avanzaba el desarrollo y se identificaban nuevas necesidades relacionadas con la generación de modelos de características. Tampoco se aplicó una metodología Scrum completa, puesto que el desarrollo no se organizó mediante sprints, roles o hitos propios de este marco de trabajo.


\subsubsection*{Control de versiones}

El control de versiones del proyecto se ha realizado mediante Git, utilizando
GitHub como plataforma de almacenamiento y seguimiento de la evolución del
código. Debido a la evolución del proyecto y a la necesidad de integrar el
generador dentro de diferentes versiones de UVLHub, el desarrollo quedó
distribuido entre varios repositorios y una solicitud de incorporación de
cambios (\textit{pull request}) en el repositorio oficial del proyecto.

Por este motivo, el repositorio
\href{https://github.com/antjimort/uvlhub}{\texttt{antjimort/uvlhub}}
no representa por sí solo todo el historial del Trabajo Fin de Grado. Dicho repositorio recoge principalmente la fase final de integración,
refactorización, corrección y validación realizada entre mayo y septiembre de
2026. La actividad desarrollada entre mayo y julio corresponde a la fase
principal de estabilización, mientras que entre el 14 y el 18 de septiembre se
realizaron correcciones adicionales derivadas de la revisión técnica de los tutores. Véase la Tabla~\ref{tab:trazabilidad-repositorios} sobre la trazabilidad de los repositorios usados durante el desarrollo.

\begin{table}[H]
\centering
\small
\renewcommand{\arraystretch}{1.15}

\begin{tabularx}{\textwidth}{@{}
    >{\raggedright\arraybackslash}p{0.25\textwidth}
    >{\centering\arraybackslash}p{0.19\textwidth}
    >{\raggedright\arraybackslash}X
@{}}
\toprule

\textbf{Recurso} &
\textbf{Periodo relevante} &
\textbf{Papel dentro del desarrollo} \\

\midrule

\href{https://github.com/antjimort/fm_generator}
{\texttt{fm\_generator}} &
15/05/2025--21/04/2026 &
Repositorio inicial del componente generador. Contiene 38 commits realizados
por el autor relacionados con la implementación de niveles UVL, atributos,
restricciones, cardinalidades y opciones de configuración. Los commits
anteriores a mayo de 2025 no se consideran parte de la contribución del TFG. \\[0.4em]

\href{https://github.com/antjimort/uvlhub_v2}
{\texttt{uvlhub\_v2}} &
14/02/2026--03/04/2026 &
Repositorio transitorio utilizado durante la migración de los cambios previos
a una nueva versión de UVLHub. En él se realizaron ajustes de migración,
persistencia de parámetros del asistente y avances en los niveles booleano y
aritmético del generador. \\[0.4em]

\href{https://github.com/diverso-lab/uvlhub/pull/117}
{\texttt{PR \#117}} &
06/04/2026--19/04/2026 &
Solicitud de incorporación de cambios en el repositorio oficial de UVLHub,
titulada \textit{feat: boolean level ready}. La PR contiene 46 commits en
total, de los cuales 45 fueron realizados por el autor, y fue fusionada en la
rama \texttt{generator} el 19 de abril de 2026. \\[0.4em]

\href{https://github.com/antjimort/uvlhub}
{\texttt{uvlhub}} &
05/2026--09/2026 &
Repositorio utilizado durante la fase final de integración, estabilización y
consolidación técnica. Incluye refactorizaciones de la arquitectura,
correcciones de compatibilidad, adaptación de pruebas, automatización de
despliegue y correcciones finales relacionadas con la validación de modelos
UVL, PySAT y la integración del generador como \textit{plugin} de Flamapy. \\

\bottomrule
\end{tabularx}

\caption{Trazabilidad de los repositorios y recursos utilizados durante el desarrollo.}
\label{tab:trazabilidad-repositorios}
\end{table}

Los repositorios y la solicitud de incorporación de cambios anteriores no
deben interpretarse como desarrollos independientes que puedan sumarse entre
sí, ya que parte de los commits realizados en el repositorio
\texttt{fm\_generator} fueron posteriormente incorporados en el PR~\#117.
La Tabla~\ref{tab:trazabilidad-repositorios} permite identificar el origen de
cada fase del desarrollo sin duplicar la dedicación contabilizada en el
cronograma.

La convención \textit{Conventional Commits} \cite{conventionalcommits} se adoptó progresivamente durante
el proyecto. Por este motivo, algunos commits iniciales no siguen una
nomenclatura completamente normalizada. Sin embargo, en las fases posteriores
se utilizaron prefijos como \textit{feat} para la incorporación de nuevas
funcionalidades, \textit{fix} para la corrección de errores, \textit{refactor}
para reorganizaciones internas del código, \textit{test} para cambios en las
pruebas, \textit{chore} para tareas de mantenimiento, \textit{style} para
ajustes de estilo y \textit{ci} para modificaciones de integración continua.

Esta organización permite mantener la trazabilidad de la evolución completa
del generador, desde la implementación inicial como componente independiente
hasta su integración y estabilización dentro de UVLHub.

\subsubsection*{Herramientas de seguimiento}

El seguimiento del desarrollo se ha realizado principalmente mediante el propio historial de GitHub, utilizando los commits realizados como registro de la evolución del proyecto y de las decisiones técnicas tomadas durante la implementación.

Además, las reuniones periódicas con los tutores han permitido revisar el estado del desarrollo para analizar las deficiencias encontradas con el propósito de corregirlas y mejorar, por tanto, la calidad del generador. Los detalles de estas reuniones se detallarán en el siguiente apartado.


\subsection{Metodología de validación y evaluación}
\label{sec:metodologia_validacion}



La validación del proyecto se ha realizado mediante un proceso iterativo durante el desarrollo, combinando pruebas funcionales, pruebas automatizadas y revisiones periódicas con los tutores del proyecto.

Tras la implementación de cada nueva funcionalidad, se realizaron comprobaciones manuales para verificar que el generador producía modelos de características correctos según los parámetros configurados, incluyendo los diferentes niveles de expresividad soportados por UVL, así como la correcta interacción del usuario con el asistente de
generación de UVLHub. Además, se ejecutaron pruebas automatizadas para
detectar errores y regresiones, incorporando casos que verificasen el
comportamiento esperado o reprodujesen los fallos relevantes detectados
durante el desarrollo.

La evaluación del cumplimiento del proyecto se ha realizado siguiendo los criterios de éxito definidos en la Tabla~\ref{tab:criterios-exito}. Estos criterios permiten comprobar de forma objetiva aspectos como la generación automática de modelos UVL, la configuración del proceso de generación, el soporte de construcciones avanzadas del lenguaje, la integración dentro del ecosistema Flamapy y UVLHub, y la existencia de una batería de pruebas automatizadas.

Además de la validación técnica continua, se realizaron revisiones del estado del proyecto junto con los tutores. Cuando se alcanzó aproximadamente el 80\% de desarrollo, se realizó una primera revisión en la que se concluyó que la solución era funcionalmente válida, pero presentaba ciertas deficiencias técnicas que debían ser corregidas antes de considerarla finalizada. Entre ellas se identificaron problemas relacionados con la integración del generador dentro de UVLHub, la pérdida de algunas funcionalidades existentes durante determinadas modificaciones (que posteriormente se recuperaron sin problema), la reducción de la cobertura de pruebas, la presencia de elementos no deseados en el repositorio y aspectos de calidad del código (documentación y limpieza de la implementación).

Tras aplicar las primeras correcciones, se realizaron nuevas revisiones
con los tutores. La versión principal obtuvo una valoración positiva en
julio de 2026. Una revisión técnica posterior identificó ajustes
adicionales relacionados con la integración con Flamapy, la validez de
los modelos generados y las pruebas; estas observaciones dieron lugar a
las correcciones efectuadas durante septiembre de 2026.

\section{Estructura de descomposición del trabajo (EDT)}
\label{sec:edt}


La Tabla~\ref{tab:edt} presenta los paquetes de trabajo y sus
entregables verificables.

\begin{longtable}{@{}P{1cm} P{4.5cm} P{8.8cm}@{}}
  \toprule
  \textbf{ID} & \textbf{Paquete de trabajo} & \textbf{Entregable verificable} \\
  \midrule
  \endhead

  1 & Gestión y planificación del proyecto &
  Planificación inicial del proyecto, estudio de la solución propuesta por los investigadores
  a partir del trabajo de referencia y seguimiento del desarrollo realizado. \\[0.3em]

  1.1 & Planificación inicial y estudio de la solución &
  Acta de constitución, definición del alcance, planificación temporal del proyecto y
  análisis de la solución propuesta en el artículo de referencia. \\[0.3em]

  1.2 & Seguimiento y control &
  Reuniones de seguimiento, revisión del estado del desarrollo y gestión de los cambios
  realizados durante el proyecto. \\[0.3em]

  2 & Análisis y diseño de la solución &
  Diseño de la arquitectura del generador y definición de los componentes necesarios para su implementación. \\[0.3em]

  2.1 & Análisis de requisitos &
  Catálogo de requisitos funcionales y técnicos del generador automático de modelos. \\[0.3em]

  2.2 & Diseño del generador y su integración &
  Arquitectura definida para la generación de modelos, configuración del sistema e integración con UVLHub. \\[0.3em]

  3 & Desarrollo del generador automático &
  Implementación del generador capaz de producir modelos de características en formato UVL. \\[0.3em]

  3.1 & Implementación de niveles de expresividad &
  Soporte implementado para generación de modelos booleanos, aritméticos y con tipos definidos por UVL. \\[0.3em]

  3.2 & Configuración e integración con el asistente &
  Sistema de configuración del generador y adaptación de su integración
  al asistente existente de UVLHub. \\[0.3em]

  4 & Pruebas y validación &
  Suite de pruebas automatizadas y validación funcional del generador integrado. \\[0.3em]

  5 & Documentación y defensa &
  Memoria final del Trabajo Fin de Grado y material necesario para la presentación de defensa. \\

  \bottomrule

  \caption{Estructura de descomposición del trabajo (EDT).}
  \label{tab:edt}
\end{longtable}

\section{Cronograma e hitos (línea base)}
\label{sec:fases}


Teniendo como fecha de inicio del proyecto el 12 de abril de 2025, el
cronograma real quedó dividido en varias etapas de desarrollo, integración y
validación. La Tabla~\ref{tab:cronograma-ejecucion} recoge los periodos de
actividad identificables a partir del historial de los repositorios y de los
hitos de integración del proyecto.

\begin{table}[H]
\centering
\small
\setlength{\tabcolsep}{4pt}
\renewcommand{\arraystretch}{1.2}

\begin{tabularx}{\textwidth}{@{}
    >{\centering\arraybackslash}p{2.0cm}
    >{\centering\arraybackslash}p{1.5cm}
    >{\centering\arraybackslash}p{1.0cm}
    >{\raggedright\arraybackslash}X
    >{\raggedright\arraybackslash}X
@{}}
\toprule

\shortstack[c]{\textbf{Periodo}\\\textbf{real}} &
\shortstack[c]{\textbf{Fin}\\\textbf{prevista}} &
\textbf{EDT} &
\textbf{Actividad principal} &
\textbf{Hito o resultado} \\

\midrule

04/2025 &
04/2025 &
1.1 &
Análisis inicial y planificación del proyecto. &
Planificación inicial y estudio de viabilidad del generador. \\[0.3em]

04/2025 &
04/2025 &
\shortstack[c]{2.1--\\2.2} &
Análisis de requisitos y diseño de la solución. &
Requisitos definidos y prototipo inicial preparado para validar la integración. \\[0.3em]

\shortstack[c]{05--07/\\2025} &
05/2026 &
3 &
Primera iteración del generador y de su integración inicial con UVLHub. &
Implementación de restricciones booleanas, cardinalidades, atributos manuales,
configuración inicial del asistente e integración preliminar del componente. \\[0.3em]

\shortstack[c]{08/2025--\\01/2026} &
--- &
--- &
Periodo de inactividad del proyecto. &
No se realizaron actividades ni se establecieron hitos durante este intervalo. \\[0.3em]

\shortstack[c]{02--04/\\2026} &
05/2026 &
3 &
Reanudación del desarrollo, migración a UVLHub 2.0 e incorporación de nuevos
niveles de expresividad. &
Migración funcional, soporte del nivel booleano, funciones agregadas e
integración mediante el PR~\#117, fusionado el 19 de abril de 2026. \\[0.3em]

\shortstack[c]{05--07/\\2026} &
05/2026 &
3--4 &
Integración final, refactorización, adaptación de pruebas y validación del
sistema. &
Generador estabilizado, pruebas adaptadas a la nueva arquitectura y validación
final de los tutores. \\

\shortstack[c]{14--26/09/\\2026} &
--- &
3--4 &
Correcciones técnicas y ampliación de las comprobaciones posteriores a
la fase principal de validación. &
Corrección de la generación de tipos y expresiones, validación de modelos UVL
con \texttt{UVLReader} y PySAT, integración del generador como \textit{plugin}
de Flamapy, reorganización de las pruebas unitarias y limpieza de elementos no
deseados del repositorio. \\

\shortstack[c]{08--10/\\2026} &
--- &
5 &
Redacción, revisión y preparación de la memoria y sus apéndices. &
Documento preparado para su entrega y defensa. \\

\bottomrule
\end{tabularx}

\caption{Cronograma de ejecución e hitos del proyecto.}
\label{tab:cronograma-ejecucion}
\end{table}

La distribución del esfuerzo se mantiene agrupada por paquetes de trabajo,
ya que las horas reales se registraron por fases generales y no de forma
independiente para cada una de las iteraciones temporales anteriores. Esta
decisión evita asignar artificialmente horas concretas a periodos para los que
no existe un registro horario detallado.

Las fases de desarrollo y su esfuerzo asociado se recogen en la
Tabla~\ref{tab:fases-desarrollo}.

\begin{table}[H]
\centering
\small
\setlength{\tabcolsep}{3pt}
\renewcommand{\arraystretch}{1.2}

\begin{tabularx}{\textwidth}{@{}
    >{\centering\arraybackslash}p{0.55cm}
    >{\raggedright\arraybackslash}X
    >{\centering\arraybackslash}p{1.35cm}
    >{\centering\arraybackslash}p{1.25cm}
    >{\centering\arraybackslash}p{1.25cm}
    >{\centering\arraybackslash}p{1.25cm}
    >{\raggedright\arraybackslash}X
@{}}
\toprule

\textbf{\#} &
\textbf{Fase} &
\textbf{EDT} &
\textbf{Obj.} &
\shortstack[c]{\textbf{Est.}\\\textbf{(h)}} &
\shortstack[c]{\textbf{Real}\\\textbf{(h)}} &
\textbf{Hito / entregable} \\

\midrule

1 &
Análisis inicial y planificación del proyecto &
1.1 &
O1 &
12 &
12 &
Prototipo inicial y planificación del desarrollo validada. \\[0.3em]

2 &
Análisis y diseño de la solución &
2.1--2.2 &
O1--O2 &
25 &
25 &
Requisitos definidos, análisis de la solución propuesta y prototipo inicial
desarrollado para validar la viabilidad de la integración. \\[0.3em]

3 &
Desarrollo e integración del generador automático &
3 &
O1--O4 &
593 &
703 &
Generador implementado e integrado dentro de UVLHub con soporte para los
niveles de expresividad definidos y corrección de errores detectados durante
el desarrollo. \\[0.3em]

4 &
Pruebas y validación de la solución &
4 &
O1, O4 y O5 &
120 &
170 &
Suite de pruebas automatizadas, validaciones funcionales y revisiones finales
del sistema completadas. \\

\midrule

&
\textbf{Total} &
&
&
\textbf{750} &
\textbf{910} &
\\

\bottomrule
\end{tabularx}

\caption{Fases de desarrollo y esfuerzo asociado: línea base frente a valores reales.}
\label{tab:fases-desarrollo}
\end{table}

La fase principal de desarrollo e integración concluyó en julio de 2026.
Posteriormente, durante septiembre, se realizaron correcciones y pruebas
adicionales derivadas de la revisión técnica de los tutores. Estas
actividades continuaron, al menos, hasta el 26 de septiembre de 2026,
fecha de la ejecución completa de las pruebas del módulo generador
documentada en este trabajo.

La línea base preveía finalizar las fases principales en mayo de 2026.
La finalización de la fase principal en julio supuso una desviación de
dos meses. La continuación de las correcciones técnicas durante septiembre
amplió a cuatro meses la desviación respecto a esa previsión.

En términos de esfuerzo, las 910 horas reales de desarrollo superaron en
160 horas las 750 horas inicialmente estimadas, lo que representa una
desviación del 21,3\,\%. El esfuerzo adicional se concentró en las fases de
desarrollo e integración del generador y de pruebas y validación, debido a
los problemas de compatibilidad con dependencias externas, la necesidad de
realizar refactorizaciones, la adaptación de las pruebas y las correcciones
finales realizadas en septiembre tras la revisión de los tutores.

De las 910 horas de desarrollo, 400 corresponden al contrato de apoyo
técnico asociado a TASOVA Plus y las 510 restantes son atribuibles al
desarrollo del TFG. Junto con las 70 horas dedicadas a la redacción,
maquetación y revisión de la memoria, las cuales se pueden ver en la Tabla~\ref{tab:fases-redaccion-memoria}, la dedicación atribuible al TFG
asciende a 580 horas.



\begin{longtable}{@{}P{1cm} P{4cm} c c c P{5cm}@{}}
\toprule
\textbf{\#} & \textbf{Capítulos} & \textbf{EDT} & \textbf{Est.} & \textbf{Real} & \textbf{Hito / Entregable} \\
\midrule
\endhead

1 &
Caps. 1--3: introducción, estado del arte y planificación
& 5
& 30 h
& 33 h
& Borrador de contexto y plan aprobado por los tutores. \\[0.3em]

2 &
Caps. 4--7: núcleo técnico (requisitos, diseño, implementación y pruebas)
& 5
& 25 h
& 24 h
& Núcleo técnico redactado y trazable con el repositorio. \\[0.3em]

3 &
Caps. 8--9: resultados y conclusiones
& 5
& 15 h
& 13 h
& Memoria completa lista para revisión final y depósito. \\

\midrule
\textbf{Total}
&
&
&
\textbf{70 h}
&
\textbf{70 h}
&
\\

\bottomrule
\caption{Fases de redacción de la memoria (estimadas vs. reales).}
\label{tab:fases-redaccion-memoria}
\end{longtable}

Las 70 horas reales indicadas en la
Tabla~\ref{tab:fases-redaccion-memoria} corresponden al trabajo completo de
elaboración de la memoria, incluyendo la redacción y revisión de los capítulos,
la organización de referencias, la preparación de tablas y figuras, la
maquetación en \LaTeX, la revisión de referencias cruzadas y la preparación de
los apéndices.

La distribución de 33, 24 y 13 horas corresponde a bloques de capítulos y no
a una estimación basada únicamente en su extensión. Las tareas técnicas
necesarias para obtener resultados, ejecutar pruebas y recopilar evidencias se
contabilizaron dentro de las fases de desarrollo y validación, evitando sumar
dos veces el mismo esfuerzo.

En la Tabla~\ref{tab:resumen-global-proyecto} se puede encontrar el resumen global en horas del proyecto.

\begin{table}[H]
\centering
\begin{tabular}{@{}P{4cm} c c@{}}
\toprule
\textbf{\# Fase} & \textbf{Est.} & \textbf{Real} \\
\midrule
Desarrollo (Tabla~\ref{tab:fases-desarrollo})
& 750 h & 910 h \\[0.3em]

Memoria (Tabla~\ref{tab:fases-redaccion-memoria})
& 70 h & 70 h \\

\midrule
\textbf{Total}
& \textbf{820 h} & \textbf{980 h} \\

\bottomrule
\end{tabular}
\caption{Resumen global del proyecto (estimadas vs. reales).}
\label{tab:resumen-global-proyecto}
\end{table}

El total de 980 horas representa la dedicación real acumulada durante el
proyecto, pero no constituye íntegramente la carga académica del TFG. De este
total, 400 horas corresponden al contrato de apoyo técnico asociado a TASOVA
Plus y no forman parte de la dedicación académica del TFG. Las 580 horas
restantes sí son atribuibles al TFG, incluyendo 510 horas de desarrollo y
70 horas dedicadas a la redacción, maquetación y revisión de la memoria.

Como referencia de carga académica, este TFG tiene 12 créditos del Sistema Europeo de Transferencia y Acumulación de Créditos (ECTS), lo
que equivale a una estimación de entre 300 y 360 horas de trabajo del
estudiante \cite{grado_software_us,rd1125_2003}. La comparación debe
realizarse con las 580 horas atribuibles al TFG, sin incluir las 400
horas correspondientes al contrato de apoyo técnico. La dedicación
registrada supera esa referencia debido a la integración con UVLHub y
Flamapy, la adaptación de las pruebas y las correcciones derivadas de
las revisiones técnicas.


\section{Plan de comunicaciones y seguimiento}
\label{sec:comunicaciones}

La Tabla~\ref{tab:comunicaciones} resume los canales utilizados para
el seguimiento del proyecto y la comunicación con los tutores.

\begin{longtable}{@{}P{3.8cm} P{2.2cm} P{2.8cm} P{5.4cm}@{}}
\toprule
\textbf{Actividad} & \textbf{Cadencia} & \textbf{Vía} & \textbf{Resultado esperado} \\
\midrule
\endhead

Reuniones con los tutores &
Según avance del proyecto &
Presencial &
Revisión del estado del desarrollo, validación de avances, detección de posibles deficiencias y definición de las próximas tareas a realizar. \\[0.3em]

Mensajería asíncrona &
Bajo demanda &
Telegram y correo electrónico &
Resolución de dudas sobre el diseño y desarrollo del sistema, recepción de indicaciones adicionales, concertación de tutorías y evaluación de avances menores. \\[0.3em]

Seguimiento del desarrollo &
Continua &
Repositorio Git/GitHub &
Trazabilidad de la evolución del proyecto mediante el historial de cambios y commits realizados durante la implementación. \\

\bottomrule
\caption{Plan de comunicaciones y seguimiento del proyecto.}
\label{tab:comunicaciones}
\end{longtable}

Durante el desarrollo se mantuvieron reuniones presenciales con los
tutores para revisar el estado de la solución y validar las decisiones
tomadas. La primera revisión relevante tuvo lugar el 5 de mayo de 2026,
cuando se había alcanzado aproximadamente el 80\,\% del desarrollo. En
ella se consideró que la solución era funcionalmente aceptable, aunque
presentaba deficiencias técnicas que debían corregirse.

Tras aplicar las primeras correcciones, el 29 de junio de 2026 se
realizó una segunda revisión y el 10 de julio se valoró positivamente
la versión principal. Posteriormente, la revisión técnica de los tutores
identificó nuevos aspectos relativos a la integración con Flamapy, la
validación de los modelos generados y la organización de las pruebas.
Estas observaciones motivaron correcciones que continuaron durante
septiembre de 2026.

\section{Gestión de riesgos}
\label{sec:riesgos}


La Tabla~\ref{tab:riesgos} recoge los riesgos identificados al inicio
del desarrollo. La exposición se ha valorado cualitativamente a partir
de la probabilidad y el impacto: se considera alta cuando el impacto es
alto y la probabilidad al menos media, o cuando la probabilidad es alta
y el impacto al menos medio.

\begin{longtable}{@{}P{0.7cm} P{3.2cm} P{1.4cm} P{1.8cm} P{2.4cm} P{4.4cm}@{}}
  \toprule
  \textbf{ID} & \textbf{Riesgo} & \textbf{Prob.} & \textbf{Impacto} & \textbf{Exposición} & \textbf{Mitigación} \\
  \midrule
  \endhead

  R1 &
  Complejidad de integración del generador dentro del ecosistema UVLHub-Flamapy &
  Alta &
  Alto &
  Alta &
  Realizar un estudio inicial de la arquitectura existente, validar progresivamente los cambios realizados y comprobar la compatibilidad entre los diferentes componentes del sistema. \\[0.3em]

  R2 &
  Problemas derivados de dependencias externas y configuración del entorno de ejecución &
  Media &
  Alto &
  Alta &
  Mantener actualizadas las dependencias necesarias, documentar la configuración del entorno y realizar pruebas de integración tras modificar versiones o paquetes externos. \\[0.3em]

  R3 &
  Errores en la generación de modelos tras la incorporación de nuevas funcionalidades &
  Alta &
  Medio &
  Alta &
  Realizar validaciones funcionales tras cada modificación e incorporar pruebas automatizadas que permitan detectar errores en las funcionalidades implementadas. \\[0.3em]

  R4 &
  Subestimación del esfuerzo necesario para completar la implementación y validación del sistema &
  Media &
  Alto &
  Alta &
  Realizar un seguimiento periódico del progreso, priorizar las tareas críticas y reservar tiempo adicional para corrección de errores y mejora de las pruebas. \\

  \bottomrule
  \caption{Riesgos identificados, exposición y planes de mitigación.}
  \label{tab:riesgos}
\end{longtable}

\subsection{Seguimiento de riesgos al cierre}
\label{sec:seguimiento_riesgos}

La Tabla~\ref{tab:seguimiento_riesgos} indica cuáles de esos riesgos
se materializaron, qué respuesta se aplicó y qué limitaciones
permanecieron al finalizar la evaluación.

\begin{longtable}{@{}P{0.7cm} P{2.8cm} P{5.4cm} P{5.3cm}@{}}
  \toprule
  \textbf{ID} & \textbf{¿Se materializó?} & \textbf{Respuesta aplicada} & \textbf{Riesgo residual} \\
  \midrule
  \endhead

  R1 &
  Sí &
  El riesgo apareció principalmente durante la fase inicial de comprensión de Flamapy y durante la integración del generador dentro de UVLHub. Para solucionarlo se estudió la arquitectura de Flamapy, se realizaron modificaciones sobre la estructura del código y se aplicó una estrategia de pruebas progresivas para validar cada incorporación realizada. &
  La integración se verificó mediante las pruebas de ejecución,
  serialización y generación desde UVLHub descritas en el
  Capítulo~\ref{chap:pruebas}. Futuras modificaciones de la arquitectura
  pueden requerir nuevas comprobaciones. \\[0.3em]

  R2 &
  Sí &
  Se produjeron problemas relacionados con las dependencias externas del proyecto, especialmente con la actualización de Flamapy, la creación y configuración de los wheels necesarios para Pyodide~\cite{pyodide} y los problemas encontrados para que estos paquetes fueran localizados y utilizados correctamente. La solución consistió en actualizar dependencias, crear nuevos wheels cuando fue necesario e investigar la configuración requerida para su correcta integración. &
  El funcionamiento se comprobó con las versiones de dependencias
  utilizadas durante la evaluación. Sus futuras actualizaciones pueden
  requerir ajustes de compatibilidad. \\[0.3em]

  R3 &
  Sí &
  Durante el desarrollo se detectaron en varias ocasiones problemas en la generación de modelos debido a que determinadas modificaciones no producían los resultados esperados. Para solucionarlos se comprobaba inicialmente el fallo mediante la generación de diferentes modelos, posteriormente se analizaba la causa, se corregía la implementación y se realizaban nuevas validaciones manuales, actualizando las pruebas unitarias cuando aportaban valor. &
  La generación se verificó mediante las pruebas automatizadas y los
  escenarios de evaluación descritos en los
  Capítulos~\ref{chap:pruebas} y~\ref{chap:resultados}. Pueden existir
  casos límite no cubiertos por esos escenarios. \\[0.3em]

  R4 &
  Sí &
  La estimación inicial del esfuerzo necesario se vio superada debido a diversos factores relacionados con la calidad final del producto, como los errores detectados durante la implementación, la necesidad de mejorar el enfoque de las pruebas, los problemas de integración y las tareas de refactorización necesarias. La gestión de este riesgo se realizó priorizando las correcciones necesarias, reorganizando tareas y utilizando las revisiones con los tutores para obtener una valoración sobre el estado del proyecto. &
  El principal riesgo residual está relacionado con la posible existencia de deuda técnica pendiente. Una mayor reducción de esta deuda mediante refactorizaciones y mejoras arquitectónicas podría facilitar futuras ampliaciones y reducir el esfuerzo necesario para incorporarlas. \\

  \bottomrule
  \caption{Seguimiento de riesgos al cierre: materialización, respuesta y residuales.}
  \label{tab:seguimiento_riesgos}
\end{longtable}

\section{Control de cambios}
\label{sec:control_cambios}

La Tabla~\ref{tab:control_cambios} registra los cambios principales
introducidos durante el desarrollo y su relación con las revisiones
del proyecto.

\begin{longtable}{@{}P{0.8cm} P{2.5cm} P{4.0cm} P{\dimexpr\textwidth-7.3cm-6\tabcolsep\relax}@{}}
\toprule
\textbf{ID} & \textbf{Fecha} & \textbf{Cambio} & \textbf{Justificación y resultado} \\
\midrule
\endhead

C1 &
12/05/2025 &
Ampliación de las opciones de configuración del generador para aumentar la flexibilidad del proceso de generación. &
\textbf{Causa:} Necesidad de proporcionar una mayor capacidad de adaptación de los modelos generados a distintos escenarios de uso.\par\textbf{Impacto:} Incremento del alcance funcional mediante la incorporación de nuevas opciones de configuración relacionadas con la generación de operadores y atributos.\par\textbf{Aprobado por:} Tutores / investigadores \\

C2 &
Durante el desarrollo &
Ampliación del plazo previsto para completar la implementación y validación del generador automático. &
\textbf{Causa:} La estimación inicial no contemplaba completamente la complejidad de la integración con el ecosistema UVLHub-Flamapy, así como los problemas surgidos durante la implementación y corrección de errores.\par\textbf{Impacto:} Desplazamiento del hito previsto de finalización del desarrollo y aumento del esfuerzo dedicado a la implementación y validación del sistema.\par\textbf{Aprobado por:} Tutores \\

C3 &
05/05/2026 &
Revisión del enfoque técnico y refactorización de la estructura interna del generador para mejorar su integración y mantenibilidad. &
\textbf{Causa:} Durante la revisión del estado del proyecto se identificaron deficiencias técnicas relacionadas con la organización interna del código y la integración del generador dentro del ecosistema existente.\par\textbf{Impacto:} Necesidad de reorganizar componentes internos, mejorar la arquitectura de la solución y adaptar parte de la implementación existente.\par\textbf{Aprobado por:} Tutores, reunión de seguimiento 05/05/2026 \\

C4 &
05/2026 &
Adaptación de la integración del generador al nuevo asistente de UVLHub,
desarrollado previamente por David Romero. &
\textbf{Causa:} La evolución del asistente existente modificó el flujo
de configuración de la generación.\par
\textbf{Impacto:} Adaptación del intercambio de parámetros y de la
ejecución del generador al asistente de seis pasos.\par
\textbf{Seguimiento:} Revisado con los tutores \\

C5 &
14--26/09/2026 &
Correcciones técnicas del generador tras la revisión de los tutores. &
\textbf{Causa:} Observaciones relativas a la integración con Flamapy,
la validez de los modelos UVL, la comprobación de satisfacibilidad y
la organización de las pruebas.\par
\textbf{Impacto:} Continuación de las correcciones y validaciones
técnicas durante septiembre de 2026, sin modificar el alcance funcional
del proyecto.\par
\textbf{Seguimiento:} Revisado con los tutores \\

\bottomrule
\caption{Registro de control de cambios sobre la línea base.}
  \label{tab:control_cambios}
\end{longtable}

\section{Plan de calidad}
\label{sec:plan_calidad}



El plan de calidad definido para el proyecto establece los criterios utilizados para garantizar la calidad tanto de la solución software desarrollada como de la documentación asociada al TFG.

\begin{itemize}

\item \textbf{Código:} se ha buscado mantener una implementación organizada,
legible y mantenible, siguiendo una estructura coherente con la arquitectura
del generador. Durante el desarrollo se realizaron tareas de refactorización
orientadas a mejorar la calidad interna del código. La trazabilidad de los
cambios se mantuvo mediante Git y el uso progresivo de la convención
\textit{Conventional Commits}. Los workflows de GitHub Actions mencionados en
la Sección~\ref{sec:ci} pertenecen al repositorio upstream de UVLHub.

\item \textbf{Entregables:} cada bloque funcional desarrollado debía cumplir con los objetivos establecidos y aportar una funcionalidad verificable dentro del sistema. Para ello, las nuevas incorporaciones se validaban progresivamente antes de continuar con el desarrollo de nuevas capacidades, asegurando que los componentes implementados mantenían el comportamiento esperado.

\item \textbf{Pruebas:} la calidad funcional del sistema se ha comprobado mediante la combinación de validaciones manuales y pruebas automatizadas. Tras la incorporación de nuevas funcionalidades se verificaba que los modelos generados cumplían los criterios esperados y que la integración con el asistente de UVLHub funcionaba correctamente. Además, se incorporaron pruebas unitarias adicionales cuando estas aportaban valor al detectar posibles errores relevantes en el sistema.

\item \textbf{Memoria:} la calidad de la documentación se ha garantizado mediante revisiones periódicas de la estructura y contenido de la memoria, así como mediante la revisión de los avances con los tutores del proyecto.

\end{itemize}

\section{Presupuesto}
\label{sec:presupuesto}



\subsection*{CapEx --- Costes de capital}

La Tabla~\ref{tab:capex} recoge los desembolsos de capital registrados.
No se realizaron compras de equipos ni pagos por licencias atribuibles
a este proyecto.

\begin{longtable}{@{}P{3.6cm} P{7cm} P{2.4cm}@{}}
\toprule
\textbf{Recurso} & \textbf{Descripción} &
\textbf{Coste directo registrado} \\
\midrule

Equipo de desarrollo &
Equipo personal ya disponible, utilizado para la implementación y las
pruebas del sistema. &
0,00\,€ \\[0.3em]

Licencias software &
Herramientas y tecnologías utilizadas sin pagos por adquisición de
licencias atribuibles al proyecto. &
0,00\,€ \\

\midrule
\textbf{Total CapEx} & & \textbf{0,00\,€} \\
\bottomrule
\caption{Desembolsos de capital registrados durante el proyecto.}
\label{tab:capex}
\end{longtable}

Los recursos utilizados durante el desarrollo no supusieron una inversión económica directa conocida, ya que se empleó principalmente infraestructura propia y herramientas basadas en software libre o de coste nulo. Por tanto, el coste real asociado a los costes de capital del proyecto es de 0,00\,€.

\subsection*{OpEx --- Costes operacionales}


\subsubsection*{Personal}


El coste de personal corresponde a la dedicación asociada al contrato de
apoyo técnico vinculado al proyecto TASOVA Plus. Dicho contrato tuvo una
duración de cuatro meses, entre abril y julio de 2025, y contemplaba una
dedicación de 25 horas semanales junto con una remuneración bruta mensual
de 930,60\,€.

La tarifa utilizada para calcular el coste estimado se ha fijado de forma
uniforme en 9,31\,€ por hora y se aplica a las 750 horas de desarrollo
previstas inicialmente. Por este motivo, los costes estimados de la tabla
se obtienen multiplicando las horas estimadas de cada fase por dicha tarifa.

El coste real refleja únicamente la remuneración correspondiente al
contrato de TASOVA Plus. Las cuatro mensualidades abonadas durante el
periodo contractual suponen un coste real total de
\(4 \times 930,60 = 3722,40\)~€, equivalente aproximadamente a
9,306\,€ por hora para las 400 horas contractuales. Las horas realizadas
específicamente para el TFG se desarrollaron fuera del periodo financiado y
no representan un coste directo adicional de personal. Véase la Tabla~\ref{tab:coste_personal}.

\begin{table}[H]
\centering
\small
\setlength{\tabcolsep}{3pt}
\renewcommand{\arraystretch}{1.25}

\begin{tabularx}{\textwidth}{@{}
    >{\raggedright\arraybackslash}X
    >{\centering\arraybackslash}p{1.35cm}
    >{\centering\arraybackslash}p{1.45cm}
    >{\centering\arraybackslash}p{1.15cm}
    >{\centering\arraybackslash}p{1.25cm}
    >{\centering\arraybackslash}p{2.05cm}
    >{\centering\arraybackslash}p{1.80cm}
@{}}
\toprule

\textbf{Actividad} &
\shortstack[c]{\textbf{H.}\\\textbf{esti-}\\\textbf{madas}} &
\shortstack[c]{\textbf{H.}\\\textbf{TASO-}\\\textbf{VA}} &
\shortstack[c]{\textbf{H.}\\\textbf{TFG}} &
\shortstack[c]{\textbf{H.}\\\textbf{reales}} &
\shortstack[c]{\textbf{Coste}\\\textbf{estimado}\\\textbf{(€)}} &
\shortstack[c]{\textbf{Salario}\\\textbf{bruto}\\\textbf{(€)}} \\

\midrule

Análisis inicial y planificación del proyecto &
12 & 12 & 0 & 12 & 111,72 & 111,67 \\

Análisis y diseño de la solución &
25 & 25 & 0 & 25 & 232,75 & 232,65 \\

Desarrollo e integración del generador automático &
593 & 363 & 340 & 703 & 5520,83 & 3378,08 \\

Pruebas y validación de la solución &
120 & 0 & 170 & 170 & 1117,20 & 0,00 \\

\midrule

\textbf{Total personal} &
\textbf{750} &
\textbf{400} &
\textbf{510} &
\textbf{910} &
\textbf{6982,50} &
\textbf{3722,40} \\

\bottomrule
\end{tabularx}

\caption{Desglose de horas y costes de personal.}
\label{tab:coste_personal}
\end{table}

\subsubsection*{Herramientas y servicios de desarrollo}

La Tabla~\ref{tab:coste_herramientas} indica los pagos registrados por
las herramientas y los servicios utilizados durante el desarrollo.

\begin{longtable}{@{}P{3cm} P{7.4cm} P{2.6cm}@{}}
\toprule
\textbf{Herramienta} & \textbf{Uso} &
\textbf{Coste directo registrado} \\
\midrule

GitHub &
Alojamiento del código fuente y seguimiento de cambios mediante el
plan utilizado durante el proyecto. &
0,00\,€ \\[0.3em]

Render &
Despliegue de una versión accesible de UVLHub durante las revisiones
funcionales. &
0,00\,€ \\[0.3em]

Python, Flask, Flamapy y dependencias &
Tecnologías utilizadas para la implementación y ejecución del sistema. &
0,00\,€ \\

\midrule
\textbf{Total herramientas} & & \textbf{0,00\,€} \\
\bottomrule
\caption{Pagos registrados por herramientas y servicios de desarrollo.}
\label{tab:coste_herramientas}
\end{longtable}

\subsection*{Resumen de costes directos cuantificados}

La Tabla~\ref{tab:presupuesto_total} resume los costes directos
cuantificados para el proyecto.

\begin{longtable}{P{4cm} P{2.5cm} P{3cm} P{3cm}}
\toprule
\textbf{Categoría} & \textbf{Tipo} & \textbf{Coste real} &
\textbf{Valor de mercado} \\
\midrule
Infraestructura y licencias & CapEx & 0,00\,€ & 0,00\,€ \\
Personal & OpEx & 3722,40\,€ & 3722,40\,€ \\
Herramientas y servicios & OpEx & 0,00\,€ & 0,00\,€ \\
\midrule
\textbf{Total cuantificado} & &
\textbf{3722,40\,€} & \textbf{3722,40\,€} \\
\bottomrule
\caption{Costes directos cuantificados para el proyecto.}
\label{tab:presupuesto_total}
\end{longtable}

El desembolso cuantificado corresponde a los 3722,40\,€ de salario
bruto asociados al contrato de apoyo técnico de TASOVA Plus. No se han
imputado compras de equipos, licencias ni pagos por los servicios de
desarrollo utilizados. Esta cifra no representa el coste económico
total del trabajo: no incluye las posibles cotizaciones a cargo de la
entidad empleadora, la amortización del equipo propio ni una valoración
monetaria de las 580 horas atribuibles al TFG.

\section{Cierre del proyecto y lecciones aprendidas}
\label{sec:cierre}


Al finalizar el proyecto, todos los paquetes de trabajo definidos en la EDT han sido completados. Las fases planificadas se han ejecutado en su totalidad, sin necesidad de descartar ningún paquete de trabajo. Los cambios relevantes realizados durante el desarrollo han quedado registrados en la Sección~\ref{sec:control_cambios}, manteniendo la coherencia con el alcance inicial definido.

La desviación temporal global del proyecto se ha concentrado principalmente en las fases de desarrollo e integración del generador automático y pruebas y validación de la solución, debido a la complejidad añadida durante la integración con componentes externos y a la necesidad de realizar iteraciones adicionales para asegurar el correcto funcionamiento del sistema.

La Tabla~\ref{tab:lecciones-aprendidas} relaciona las principales
dificultades encontradas con los cambios que se aplicarían en un
desarrollo similar.

\begin{longtable}{P{6cm} P{6cm}}
  \toprule
  \textbf{Lección} & \textbf{Qué se haría distinto} \\
  
  \midrule
  
  La estimación inicial del esfuerzo necesario para las fases de integración y validación fue inferior al esfuerzo realmente requerido. &
  Realizar una estimación más conservadora, incluyendo un margen adicional para problemas derivados de dependencias externas y correcciones durante la fase final. \\[0.5em]
  
  El conocimiento inicial sobre algunas tecnologías utilizadas en el proyecto no era suficiente para prever completamente las dificultades de integración. &
  Dedicar una fase inicial más extensa al estudio de las herramientas y componentes externos antes de comenzar la implementación. \\[0.5em]
  
  La estrategia de pruebas podría haberse definido con mayor profundidad desde las primeras etapas del desarrollo. &
  Diseñar desde el inicio una planificación de pruebas más completa, identificando las funcionalidades críticas y los casos de validación necesarios. \\[0.5em]
  
  La acumulación progresiva de cambios durante el desarrollo generó cierta deuda técnica. &
  Realizar revisiones y refactorizaciones periódicas del código para mantener una arquitectura más mantenible durante todo el ciclo de desarrollo. \\
  
  \bottomrule
  \caption{Lecciones aprendidas al cierre del proyecto.}
  \label{tab:lecciones-aprendidas}
\end{longtable}

En conjunto, el proyecto se ha completado con una dedicación real de
980 horas, de las cuales 910 corresponden al desarrollo y 70 a la elaboración
de la memoria. Aunque la planificación inicial se vio afectada por un periodo
de inactividad y por varias iteraciones adicionales de integración y
validación, todos los paquetes de trabajo definidos en la EDT fueron
completados.

Con el proyecto planificado, ejecutado y cerrado, el siguiente capítulo
precisa qué debe hacer exactamente la solución desarrollada. El
Capítulo~\ref{chap:requisitos} presenta el catálogo de requisitos funcionales
y no funcionales que sirve como base para el diseño, la implementación y la
validación del sistema.

